% =====================================================================
%  A. Hirschowitz, C. Simpson, "Descente pour les n-champs",
%  arXiv:math/9807049v3 (2001).  English translation.
%
%  Build:  latexmk -pdf -outdir=build src/main.tex
%          (or simply: make, from this directory)
%
%  main.tex carries no page metadata of its own; every content file does.
%  Page numbers in %@ pages are those of the original PDF (1--250).
%
%  The mathematics, notation, labels and numbering of results are the
%  original's.  Only the prose is translated.  Do not "improve" a formula:
%  the translation must remain checkable against the French source.
% =====================================================================

\documentclass[11pt,a4paper]{article}

% ------------------------------------------------------------------ packages
\usepackage[T1]{fontenc}
\usepackage[utf8]{inputenc}
\usepackage{lmodern}
\usepackage[english]{babel}
\usepackage{amsmath}
\usepackage{amssymb}
\usepackage{amsthm}
\usepackage{mathtools}
\usepackage[a4paper,margin=2.6cm]{geometry}
\usepackage{xcolor}
\usepackage{hyperref}

\hypersetup{
  colorlinks=true,
  linkcolor=black!70!blue,
  citecolor=black!70!blue,
  urlcolor=black!70!blue,
  pdftitle={Descent for n-Stacks},
  pdfauthor={André Hirschowitz and Carlos Simpson},
  pdfsubject={English translation of Descente pour les n-champs, arXiv:math/9807049v3},
}

% ------------------------------------------------------------------ numbering
% The original is an article with 21 numbered sections and unnumbered
% subsection headings (\subnumero).  Sections and equations keep the
% original's numbers.
\renewcommand{\thesection}{\arabic{section}}
\setcounter{secnumdepth}{1}
\setcounter{tocdepth}{2}

% ------------------------------------------------------------------ theorems
% The original declares one counter, `lemma', numbered within the section,
% and makes theorem, proposition, corollary, definition, problem, conjecture,
% convention and remark share it:
%
%   \newtheorem{lemma}{Lemma}[section]
%   \newtheorem{theorem}[lemma]{Theorem}
%   ...
%
% Because the French and English environment names were mixed in the source,
% the original body sometimes reads "Définition 2.3" and sometimes
% "Definition 2.3", but the number is always the shared section counter.
% The translation prints every environment in English while keeping the
% exact same counter and therefore the exact same numbers.
\newtheorem{lemma}{Lemma}[section]
\newtheorem{theorem}[lemma]{Theorem}
\newtheorem{corollary}[lemma]{Corollary}
\newtheorem{proposition}[lemma]{Proposition}
\newtheorem{problem}[lemma]{Problem}
\newtheorem{conjecture}[lemma]{Conjecture}
\newtheorem{convention}[lemma]{Convention}
\newtheorem{exercise}[lemma]{Exercise}
\newtheorem{lemmedefinition}[lemma]{Lemma-Definition}

\theoremstyle{definition}
\newtheorem{definition}[lemma]{Definition}
\newtheorem{remark}[lemma]{Remark}

% French environment names, kept as aliases so that a source file that has
% not yet been converted still compiles and still prints an English label.
\newtheorem{lemme}[lemma]{Lemma}
\newtheorem{theoreme}[lemma]{Theorem}
\newtheorem{corollaire}[lemma]{Corollary}
\newtheorem{probleme}[lemma]{Problem}
\newtheorem{remarque}[lemma]{Remark}

% The original's unnumbered corollary environment (if any survive).
\newtheorem*{unnumberedcorollary}{Corollary}

% ------------------------------------------------------------------ headings
% Section headings are numbered; the original's \subnumero headings are
% unnumbered but appear in the bookmarks.
\newcommand{\numero}[1]{\clearpage\section{#1}}
\newcommand{\subnumero}[1]{%
  \subsection*{#1}%
  \phantomsection\addcontentsline{toc}{subsection}{#1}}

% Bold run-in paragraph heading, as in the original.
\newcommand{\paragraf}[1]{\vspace{.1in}\par\noindent
  {\bf #1.}\hspace*{.12in}}

% End-of-proof box.
\newcommand{\eop}{\hfill $\Box$\vspace{.1in}}

% The original numbers displayed equations with a private counter that is
% printed as (section.n).  Reproduce that exactly.
\newcommand{\equanum}{\vspace{.1in} \newline
  \addtocounter{subsection}{1}(\thesubsection )
  \vspace{-.21in}}

\newcommand{\subequanum}{\vspace{.1in} \newline
  \addtocounter{subsubsection}{1}(\thesubsubsection )
  \vspace{-.21in}}

\newcommand{\numeroOut}[2]{
  \vspace{.5in}\par\noindent\addtocounter{subsection}{1}(\thesubsection
  )\,\, {\em #2}\,\,\, {\tt (#1)}\newline}

\newenvironment{inset}{
  \begin{list}{}{\setlength{\rightmargin}{.1in}
    \setlength{\leftmargin}{.2in}\setlength{\topsep}{.08in}
    \setlength{\listparindent}{.13in}\setlength{\itemindent}{.13in}}
  \small\item}{
  \end{list}\normalsize}

\newenvironment{smallset}{
  \begin{list}{}{\setlength{\rightmargin}{0in}
    \setlength{\leftmargin}{0in}\setlength{\topsep}{.08in}
    \setlength{\listparindent}{.2in}\setlength{\itemindent}{.2in}}
  \small\item}{
  \end{list}\normalsize}

% The original's capitalised Theorem environment.
\newenvironment{Theorem}{\vspace{.05in} \par
  \noindent {\bf Theorem} \hspace*{.05in} \em }{
  \vspace{.05in}\newline }

% ------------------------------------------------------------------ notation
% Copied verbatim from the modernized source: the mathematics must not move.
\renewcommand{\boxtimes}{\mbox{\hspace*{.05
cm}$\times$\hspace*{-.33
cm}\raisebox{-.04cm}{$\Box$}}}

\newcommand{\semidirect}{\mbox{\hspace*{.12
cm}$\times$\hspace*{-.38
cm}\raisebox{0cm}{$\triangleright$}\hspace*{.28cm}}}

\newcommand{\cc}{{\bf C}}
\newcommand{\rr}{{\bf R}}
\newcommand{\qq}{{\bf Q}}
\newcommand{\hh}{{\bf H}}
\newcommand{\qbar}{\overline{\bf Q}}
\newcommand{\zz}{{\bf Z}}
\newcommand{\pp}{{\bf P}}
\newcommand{\AAA}{{\bf A}}

\newcommand{\mm}{{\bf M}}
\newcommand{\nn}{{\bf N}}

\newcommand{\Ox}{{\cal O}_X}
\newcommand{\Dx}{{\cal D}_X}

\newcommand{\Cc}{{\cal C}}
\newcommand{\Ee}{{\cal E}}
\newcommand{\Ff}{{\cal F}}
\newcommand{\Gg}{{\cal G}}
\newcommand{\Qq}{{\cal Q}}
\newcommand{\Tt}{{\cal T}}
\newcommand{\Oo}{{\cal O}}
\newcommand{\Dd}{{\cal D}}
\newcommand{\Aa}{{\cal A}}
\newcommand{\Bb}{{\cal B}}
\newcommand{\Mm}{{\cal M}}
\newcommand{\Ww}{{\cal W}}
\newcommand{\Kk}{{\cal K}}
\newcommand{\Vv}{{\cal V}}
\newcommand{\Uu}{{\cal U}}
\newcommand{\Ll}{{\cal L}}
\newcommand{\Hh}{{\cal H}}
\newcommand{\Ii}{{\cal I}}
\newcommand{\Jj}{{\cal J}}
\newcommand{\Nn}{{\cal N}}
\newcommand{\Pp}{{\cal P}}
\newcommand{\Rr}{{\cal R}}
\newcommand{\Ss}{{\cal S}}
\newcommand{\Xx}{{\cal X}}
\newcommand{\Yy}{{\cal Y}}
\newcommand{\Zz}{{\cal Z}}

\newcommand{\Hilb}{\mbox{\bf Hilb}}
\newcommand{\Grass}{\mbox{\bf Grass}}
\newcommand{\Mod}{{\bf M}}
\newcommand{\VHS}{{\bf V}}
\newcommand{\Gm}{{\bf G}_m}
\newcommand{\Rep}{{\bf R}}
\newcommand{\Def}{{\bf D}}
\newcommand{\Alg}{\mbox{\bf Alg}}
\newcommand{\Inf}{\mbox{\bf Inf}}

\newcommand{\DD}{{\cal D}^+}
\newcommand{\Dbc}{{\cal D}^b_c}

\newcommand{\Cchi}{\begin{array}{c} \rule[-.12in]{0in}{.2in}
{\bf \chi}
\end{array}\!}

\newcommand{\gera}{{\bf a}}
\newcommand{\gerb}{{\bf b}}
\newcommand{\gerc}{{\bf c}}
\newcommand{\gerd}{{\bf d}}
\newcommand{\gere}{{\bf e}}
\newcommand{\gerf}{{\bf f}}
\newcommand{\gerg}{{\bf g}}
\newcommand{\gerh}{{\bf h}}
\newcommand{\geri}{{\bf i}}
\newcommand{\gerj}{{\bf j}}
\newcommand{\gerk}{{\bf k}}
\newcommand{\gerl}{{\bf l}}
\newcommand{\germ}{{\bf m}}
\newcommand{\gern}{{\bf n}}
\newcommand{\gero}{{\bf o}}
\newcommand{\gerp}{{\bf p}}
\newcommand{\gerq}{{\bf q}}
\newcommand{\gerr}{{\bf r}}
\newcommand{\gers}{{\bf s}}
\newcommand{\gert}{{\bf t}}
\newcommand{\geru}{{\bf u}}
\newcommand{\gerv}{{\bf v}}
\newcommand{\gerw}{{\bf w}}
\newcommand{\gerx}{{\bf x}}
\newcommand{\gery}{{\bf y}}
\newcommand{\gerz}{{\bf z}}

\newcommand{\delbar}{\overline{\partial}}
\newcommand{\Fperp}{F^{\perp}}
\newcommand{\thetbar}{\overline{\theta}}

\newcommand{\dnu}{\frac{\partial}{\partial\nu}}
\newcommand{\Xphi}{X_{\varphi}}
\newcommand{\Yphi}{Y_{\varphi}}
\newcommand{\Hphi}{H_{\varphi}}
\newcommand{\Hphit}{H_{\varphi ,t}}

\newcommand{\Choose}[2]{\left(\begin{array}{c} #1 \\ #2
\end{array}\right)}

\newcommand{\bra}[1]{\langle #1 |}
\newcommand{\ket}[1]{| #1 \rangle}

\newcommand{\AutT}{Aut^{\otimes}}
\newcommand{\autT}{\gera\geru\gert ^{\otimes}}

\newcommand{\Cbar}{\overline{\cal C}}
\newcommand{\Chat}{\hat{\cal C}}
\newcommand{\dhat}{\hat{d}}

\newcommand{\CDR}{{\cal C}_{DR}}
\newcommand{\CDol}{{\cal C}_{Dol}}
\newcommand{\CHH}{{\cal C}_{H}}
\newcommand{\CDp}{{\cal C}_{D'}}

\newcommand{\gap}{{\rm gap}}

% The source's French connectors inside formulas.  Kept only so that an
% unconverted passage compiles; translators replace them by English prose.
\newcommand{\remplacant}{replacing\,\, }
\newcommand{\ca}{this\,\, }
\newcommand{\facon}{way\,\, }

\newcommand{\tworightarrows}{\stackrel{\displaystyle \rightarrow}{\rightarrow}}

% Reedy-category notation.
\newcommand{\direct}[1]{\stackrel{\rightarrow}{#1}}
\newcommand{\inverse}[1]{\stackrel{\leftarrow}{#1}}

% ------------------------------------------------------------------ editorial
% Translator's note: used for an emendation where the French prints
% something that is wrong as mathematics.  The note says what was printed
% and what we print instead.  See ../../AGENTS.md.
\newcommand{\ednote}[1]{\footnote{\emph{Translator's note.} #1}}

% Placeholder printed in place of a passage that has not been translated yet.
\newcommand{\notyet}[1]{%
  \par\noindent
  \begingroup
  \setlength{\fboxsep}{6pt}%
  \fcolorbox{black!40}{black!4}{%
    \parbox{\dimexpr\linewidth-2\fboxsep-2\fboxrule\relax}{%
      \small\itshape
      Not yet translated --- #1\par}}%
  \endgroup\par}

% Long formulas inside lists otherwise overfull.
\setlength{\emergencystretch}{3em}

% ------------------------------------------------------------------ metadata
\title{Descent for $n$-Stacks}
\author{Andr\'e Hirschowitz \and Carlos Simpson}
\date{1998; revised 2001}
